% =========================================================================
% Compute-Aware Anytime MPPI for UAV Control — ECC submission draft
% Phase 5 full draft (2026-07-18). Numbers are from the committed Phase 2
% datasets (results/ @20Hz n=13/condition, results_40hz/ @40Hz n=8/condition).
% Locked wording rules: PROJECT_PLAN.md "Paper wording (locked)".
% Build: pdflatex main && bibtex main && pdflatex main && pdflatex main
% =========================================================================
\documentclass[conference]{IEEEtran}
\usepackage{amsmath,amssymb}
\usepackage{booktabs}
\usepackage{graphicx}
\usepackage{cite}

\usepackage{xcolor}
\newcommand{\todo}[1]{\textbf{\textcolor{red}{[TODO: #1]}}}

\title{Compute-Aware Anytime MPPI: Runtime Sample-Count Adaptation\\
for Real-Time UAV Control}

% TODO: author block (name, affiliation, email; funding acks)
% Byline uses the author's middle name by preference: Ishan Jakkulwar.
\author{\IEEEauthorblockN{Ishan Jakkulwar}
\IEEEauthorblockA{\todo{affiliation / email / advisor co-authors}}}

\begin{document}
\maketitle

% =========================================================================
\begin{abstract}
Sampling-based optimal control methods such as Model Predictive Path
Integral (MPPI) control require selecting a number of trajectory samples
per control cycle. Existing UAV deployments treat this sample count as a
fixed, offline-tuned design parameter, implicitly assuming that the compute
available to the controller is constant at runtime --- an assumption that
fails when the controller shares a companion computer with perception or
other variable workloads. We apply feedback-scheduling principles to
MPPI's sampling budget: measured per-cycle execution time drives a
utilization-targeting budget law that adapts computational effort while
maintaining closed-loop control deadlines, with an explicit safety
fallback mechanism. The mechanism is deliberately classical; the
contribution is its application to a stochastic sampling-based optimizer
--- whose computation is fungible in fine increments with gracefully
degrading quality --- and a pre-registered, two-baseline, two-regime
experimental evaluation of what adaptation does and does not buy.
In software-in-the-loop experiments (ROS~2, PX4, Gazebo) across repeated
trials, the adaptive controller reduced deadline misses from
$18.0 \pm 4.1$ to $0.15 \pm 0.38$ per trial relative to a conventional
fixed-budget baseline under CPU contention ($p < 10^{-3}$), with
statistically indistinguishable tracking accuracy. Against a
\emph{carefully tuned} constant budget --- fixed a priori at the adaptive
scheduler's preliminary average --- adaptation performed equivalently in
the regime the constant was tuned for, and we show why: adaptation
operates on the flat region of MPPI's quality--compute curve. In a second,
pre-registered operating regime with a tighter control period, the tuned
constant missed 30.5\% of all control cycles while the adaptive controller
re-regulated its budget automatically and maintained essentially zero
misses ($0.75 \pm 0.71$ per trial). The value of runtime adaptation is
thus robustness across operating regimes without offline retuning, rather
than outperforming a well-tuned constant within any single regime.
\end{abstract}

% =========================================================================
\section{Introduction}
\label{sec:intro}

Model Predictive Path Integral (MPPI) control has emerged as a practical
sampling-based optimal control method for agile robotic platforms: it
handles nonlinear dynamics and non-convex, even non-smooth cost functions
by evaluating a large number $N$ of stochastic trajectory rollouts per
control cycle and averaging their controls under an information-theoretic
weighting~\cite{williams2016aggressive,williams2018information}. On
unmanned aerial vehicles (UAVs), MPPI has recently been demonstrated with
onboard, embedded computation~\cite{minarik2024mppi,enrico2025comparison}.

The per-cycle computational cost of MPPI is essentially linear in $N$,
and the choice of $N$ mediates a direct trade-off between control quality
and compute. In every embedded deployment we are aware of, this trade-off
is resolved \emph{offline}: the designer characterizes the target
hardware, selects the largest $N$ that fits the control period, and holds
it fixed~\cite{enrico2025comparison}. A fixed budget is a reasonable
default --- it simplifies timing analysis, implementation, and tuning ---
but it hard-codes an assumption that rarely holds in practice: that the
compute available to the controller is \emph{constant} at runtime. On a
real aerial platform the controller typically shares a companion computer
with perception, mapping, state estimation, logging, and communication
stacks whose demands fluctuate. When a compute transient collides with a
fixed sampling budget, the controller either misses its deadline ---
degrading the real-time contract that the surrounding autopilot stack
depends on --- or must be provisioned so conservatively that quality is
sacrificed in the common case.

This paper treats the sample count as a \emph{runtime-adaptive control
variable} instead. We formulate sample-count selection as an online
computational resource allocation problem: a lightweight scheduler
measures the controller's own recent per-sample execution time, allocates
the next cycle's budget to fit a fixed fraction of the control period, and
triggers a safety fallback under consecutive deadline misses. The
mechanism is deliberately simple --- a measured-cost budget law with
clamping and smoothing --- and requires no model of the contention source.

\subsubsection*{Contributions}
\begin{enumerate}
  \item The application of feedback-scheduling
  principles~\cite{lu2002feedback,cervin2002feedback} to the sampling
  budget of a stochastic sampling-based optimal controller: a
  measured-execution-time budget law with a deadline-satisfaction
  objective and a domain-specific safety fallback, the latter validated
  in flight by fault injection (Sec.~\ref{sec:results:fault}). We claim
  no novelty for the scheduling mechanism itself; the novelty lies in
  the actuator and its consequences.
  \item An offline characterization of MPPI's quality--compute trade-off
  (a Pareto sweep over $N$) that both motivates the approach and, as we
  show, \emph{predicts} its closed-loop behavior: adaptation operates on
  the flat region of the quality--compute curve, which is precisely why
  large budget reductions cost little tracking quality
  (Sec.~\ref{sec:results}).
  \item A pre-registered, staged experimental validation in
  software-in-the-loop (ROS~2 Humble, PX4, Gazebo Harmonic, MAVROS):
  repeated trials against \emph{two} baselines --- the conventional
  fixed-maximum budget and a \emph{predetermined, carefully tuned}
  constant budget --- across \emph{two} operating regimes (50\,ms and
  25\,ms control deadlines), with statistical tests
  (Sec.~\ref{sec:setup},~\ref{sec:results}).
  \item An honest negative result with a systems explanation: within the
  regime a constant budget is tuned for, adaptation is statistically
  indistinguishable from it --- and on our platform, OS scheduling
  shields a periodic control task strongly enough that no fair-share CPU
  load level could break the tuned constant at the 50\,ms deadline. The
  case for adaptation is robustness \emph{across} regimes: the tuned
  constant missed 30.5\% of cycles at the tighter deadline, while
  adaptation transferred with zero retuning (Sec.~\ref{sec:discussion}).
\end{enumerate}

% =========================================================================
\section{Related Work}
\label{sec:related}

Our work sits at the intersection of three threads: sampling-based MPC on
aerial platforms, computational-resource-aware control, and adaptive
sampling within MPPI itself.

\subsection{MPPI on aerial platforms and embedded hardware}

MPPI samples candidate control trajectories per cycle and combines them
via an information-theoretic
weighting~\cite{williams2016aggressive,williams2018information}. The
sample count $N$ is a central parameter, determining both the quality of
the optimized control and, linearly, the per-cycle cost. A substantial
body of work has extended MPPI along several axes: robustness via
ancillary tracking controllers and augmented-state
formulations~\cite{williams2018robust,gandhi2021robust}, improved
sampling distributions~\cite{mohamed2022logmppi}, generalized
information-theoretic objectives~\cite{wang2021tsallis}, and safety
guarantees through control barrier functions~\cite{yin2023shield}.
Existing MPPI research has thus largely focused on improving the sampling
distribution, robustness, constraint handling, or computational
acceleration; the \emph{size} of the sampling budget itself, and its
relationship to fluctuating runtime compute availability, has received
comparatively little attention.

On UAVs specifically, Mina\v{r}\'ik et al.~\cite{minarik2024mppi}
demonstrate onboard GPU-parallelized MPPI in real flight, noting that
real-time feasibility constrains the usable sample count. Enrico, Mancini
and Capello~\cite{enrico2025comparison} compare NMPC against
GPU-parallelized MPPI on a Jetson Orin Nano under a ROS~2/PX4 stack
closely matching ours, selecting an MPPI configuration (horizon 40,
800--1250 samples) explicitly to satisfy the 50\,Hz deadline imposed by
the flight controller --- a \emph{fixed} configuration derived through
\emph{offline} analysis. Across representative embedded MPPI
implementations, the sampling budget is selected offline based on
expected computational resources and then held fixed during deployment.
This choice is not accidental: a fixed budget simplifies timing analysis,
implementation, and parameter tuning, making it the de facto standard ---
but it embeds the assumption that compute availability is constant at
runtime.

\subsection{Computational-resource-aware control}

A separate literature treats computation explicitly as a limited resource
for control. \emph{Anytime control} designs controllers whose output
quality improves with allotted computation, so early termination still
yields a usable command: switching among control laws of increasing
quality under processor preemption~\cite{fontanelli2008anytime},
computing control sequences under stochastic processor
availability~\cite{quevedo2013sequence}, and --- closest in spirit to our
setting --- co-designing an anytime perception-based \emph{estimator}
with a robust controller using an offline-characterized delay/error
trade-off curve as a runtime contract~\cite{pant2021anytime}. These
frameworks established the principle we build on, but target
deterministic or linear formulations (or the estimation side of the
loop), not the internal sampling budget of a stochastic sampling-based
optimizer.

\emph{Feedback scheduling.} The real-time systems community has long
applied feedback control to computation itself: feedback control
real-time scheduling regulates CPU utilization via measured execution
times~\cite{lu2002feedback}; the elastic task model adapts task rates
within declared ranges under overload~\cite{buttazzo1998elastic}; and
feedback--feedforward scheduling applies these ideas to control tasks
specifically, adjusting sampling periods from measured
load~\cite{cervin2002feedback}. Our budget law
(Sec.~\ref{sec:scheduler}) is an instance of this family --- a
measured-cost, utilization-targeting allocator with smoothing and
clamping --- and we claim no novelty for the mechanism itself. What is
new, to the best of our knowledge, is the \emph{actuator}: rather than
adjusting task periods or admission, we adjust the internal sampling
budget of a stochastic sampling-based optimizer, exploiting a property
classical feedback scheduling could not assume --- MPPI's per-cycle
computation is fungible in fine increments, with a characterizable and
gracefully degrading quality--compute relationship
(Sec.~\ref{sec:pareto}), rather than the step change of dropping or
delaying a whole task.

\emph{Event- and self-triggered MPC} adapts \emph{when} control is
computed rather than \emph{how much} computation each solve
receives~\cite{heemels2012introduction,gommans2015resource}.
\emph{Weakly-hard real-time} models formalize tolerance to bounded
deadline-miss patterns~\cite{bernat2001weakly}, with recent work
characterizing stability under consecutive-miss
constraints~\cite{maggio2020control}. That thread treats the miss pattern
as given and analyzes its consequences; our scheduler acts
\emph{upstream}, adapting per-cycle computation to avoid misses, with the
weakly-hard perspective informing our consecutive-miss safety fallback.

\subsection{Adaptive sampling within MPPI}

A third, orthogonal thread adapts MPPI's sampling \emph{distribution}:
Asmar et al.~\cite{asmar2023mpopis} generalize MPPI to admit adaptive
importance sampling, matching proposal moments to improve sample
efficiency at fixed budget. Such methods adapt \emph{where} to sample,
driven by optimization performance; they do not adapt \emph{how many}
samples to draw, and the adaptation is not driven by measured compute
availability.

\subsection{Positioning}

We are not aware of prior work that applies feedback-scheduling
principles to the sampling budget of a stochastic sampling-based optimal
controller, with an explicit deadline-satisfaction objective and safety
fallback, validated in closed-loop UAV experiments. Our results connect
the threads above: an
offline quality--compute characterization (in the spirit
of~\cite{pant2021anytime}) predicts closed-loop quality outcomes across
operating regimes, because adaptation operates on the flat region of
MPPI's quality--compute curve, while deadline adherence --- the
weakly-hard-motivated metric --- is where adaptation shows its value.

% =========================================================================
\section{Method}
\label{sec:method}

\subsection{Preliminaries: MPPI}

We consider velocity-commanded quadrotor position control with state
$\mathbf{x}_t \in \mathbb{R}^3$ (position) and control
$\mathbf{u}_t \in \mathbb{R}^3$ (commanded velocity), propagated over a
horizon of $H$ steps with timestep $\Delta t$ under single-integrator
prediction dynamics $\mathbf{x}_{k+1} = \mathbf{x}_k + \mathbf{u}_k
\Delta t$ (the plant itself is the full PX4/Gazebo quadrotor stack; see
Sec.~\ref{sec:threats} for the model-mismatch discussion). At each control
cycle, MPPI draws $N$ control-sequence samples
$\{\mathbf{u}^{(i)}_{0:H-1}\}_{i=1}^{N}$ by perturbing a nominal sequence
with Gaussian noise of standard deviation $\sigma$, rolls each out through
the prediction model, and scores it with a quadratic cost
\begin{equation}
S^{(i)} = \sum_{k=0}^{H-1} \Big( \|\mathbf{x}^{(i)}_k -
\mathbf{x}^{\mathrm{ref}}\|^2_{Q} + \|\mathbf{u}^{(i)}_k\|^2_{R} \Big)
 + \|\mathbf{x}^{(i)}_H - \mathbf{x}^{\mathrm{ref}}\|^2_{Q_f},
\end{equation}
then combines samples by exponentiated-cost weighting with inverse
temperature $\lambda$:
\begin{equation}
w_i = \frac{\exp\!\big(-\tfrac{1}{\lambda}(S^{(i)} - S_{\min})\big)}
{\sum_{j=1}^{N}\exp\!\big(-\tfrac{1}{\lambda}(S^{(j)} - S_{\min})\big)},
\qquad
\mathbf{u}^\star = \sum_{i=1}^{N} w_i\, \mathbf{u}^{(i)}_{0}.
\end{equation}
The first control of the weighted average is published as the velocity
setpoint. The per-cycle cost of this procedure is dominated by the $N$
rollouts and is, to good approximation, linear in $N$.

\subsection{The quality--compute trade-off}
\label{sec:pareto}

We characterized quality versus $N$ offline with a 10-seed sweep over
$N \in \{20, 50, 100, 150, 200, 250, 300, 350, 400\}$ on a standalone
(simulator-free) tracking task. RMS tracking error improves by only
13.4\% from $N{=}20$ to $N{=}400$ while per-call compute grows
$\sim$19$\times$; beyond $N \approx 150$--$200$ the differences are
within the sweep's own seed-to-seed noise floor
($\sigma \approx 0.03$ in the sweep's units). The curve is thus steep
below $N \approx 150$ and essentially flat above --- an
\emph{anytime-friendly} shape: samples beyond the elbow are cheap to
give up. \todo{Figure: Pareto sweep, mean $\pm$ std vs.\ $N$.}

\subsection{Anytime scheduler}
\label{sec:scheduler}

Let $T$ be the control period (deadline), $\beta \in (0,1)$ a budget
margin, and $c_t$ the measured wall-clock duration of cycle $t$'s MPPI
call with $N_t$ samples. The scheduler maintains a smoothed per-sample
cost estimate over a window of the last $w$ cycles,
\begin{equation}
\hat{\tau}_t = \frac{1}{w}\sum_{j=t-w+1}^{t} \frac{c_j}{N_j},
\end{equation}
and allocates the next cycle's budget as
\begin{equation}
\label{eq:budget}
N_{t+1} = \mathrm{clamp}\!\left(\left\lfloor
\frac{\beta\, T}{\hat{\tau}_t} \right\rfloor,\; N_{\min},\; N_{\max}\right).
\end{equation}
Equation~\eqref{eq:budget} is a classical feedback-scheduling law ---
measured per-unit cost, utilization-target allocation, smoothing,
clamping --- in the lineage of feedback control real-time
scheduling~\cite{lu2002feedback} and elastic task
adaptation~\cite{buttazzo1998elastic}; what differs is the actuator (an
optimizer's internal sampling budget rather than task periods or
admission). It is a ratio law: it targets a compute utilization of
$\beta T$ per cycle regardless of \emph{why} per-sample cost changed
(CPU contention, frequency scaling, or a changed deadline), which is
what allows the same law to transfer across operating regimes
(Sec.~\ref{sec:results:40hz}). We use $\beta = 0.6$, $w = 3$,
$N_{\min} = 20$, $N_{\max} = 400$; the margin $1 - \beta$ absorbs
measurement noise, scheduling jitter, and the non-MPPI portion of the
cycle.

\emph{Safety fallback.} Independently of the smoothed estimate, the
scheduler counts consecutive raw deadline overruns ($c_t > T$); at two
consecutive misses it signals a fallback condition under which a simpler
controller (e.g., the PID stage of our stack) can be substituted. The
threshold of two avoids flapping on isolated timing spikes and is
motivated by consecutive-miss weakly-hard stability
analyses~\cite{maggio2020control}. Under the comparison conditions of
Sec.~\ref{sec:results} the adaptive controller never triggered the
fallback --- itself evidence of the budget law's margin adequacy --- so
we validate the mechanism by fault injection: a synthetic compute stall,
inserted inside the timed section of the control cycle and thus
indistinguishable from a real compute transient, forces consecutive
deadline overruns and exercises the trigger/recovery path in flight
(Sec.~\ref{sec:results:fault}).

\subsection{Implementation}

The controller runs in a ROS~2 Humble node at a wall-timer period $T$
(20\,Hz or 40\,Hz), consuming MAVROS pose estimates and publishing MAVROS
velocity setpoints to PX4. The node performs no control mathematics
itself; all optimization lives in controller classes, and the scheduler
interacts with MPPI solely through the measured call duration and the
per-call sample-count override --- the mechanism is agnostic to the MPPI
internals and would compose with the distribution-adaptation methods of
Sec.~\ref{sec:related}. Cycle-level telemetry (sample count, call time,
state-to-command latency, tracking error, deadline flag) is logged to CSV
every cycle.

% =========================================================================
\section{Experimental Setup}
\label{sec:setup}

\subsection{Stack, task, and conditions}

Experiments use PX4 SITL with Gazebo Harmonic (x500 quadrotor), MAVROS,
and ROS~2 Humble on a 16-hardware-thread Intel i7-13620H running Linux
6.8. The task is a 4-waypoint square circuit (2\,m sides at 5\,m
altitude) followed by station-keeping at the final waypoint; each trial
comprises takeoff, arming/offboard entry, and 90--105\,s of tracking.

Three controller conditions are compared:
\begin{itemize}
  \item \textbf{Adaptive}: the scheduler of Sec.~\ref{sec:scheduler}.
  \item \textbf{Fixed-400}: the conventional deployment choice,
  $N = N_{\max} = 400$ always.
  \item \textbf{Constant-330}: an equivalent-budget control condition.
  The value was \emph{fixed a priori}, before any Phase-2 trial ran, at
  the adaptive scheduler's average $N = 329.77$ measured in a
  preliminary stress experiment --- deriving it from the main
  experiment's own adaptive runs would make the baseline depend on the
  result it is compared against. It was not re-tuned afterwards.
\end{itemize}

\subsection{Stress condition (calibrated, then locked)}
\label{sec:stress}

Compute contention is injected by a load generator sharing the
controller's CPU set. Because a general-purpose OS scheduler (CFS/EEVDF)
strongly shields a periodic control task from fair-share busy-wait load
--- in calibration, 32 and 64 unpinned threads and 16 threads pinned to
4 cores all left the Fixed-400 call time at $\sim$37\,ms with zero misses
at the 50\,ms deadline, while a single-core pin destabilized flight
entirely --- the locked condition pins the controller node and the load
to \emph{two} cores (an embedded-class compute envelope; tighter
real-time constraints of this kind are expected on lower-performance
embedded processors) with 6 busy-wait threads for a 30\,s window
beginning 20\,s after tracking starts. The simulator, autopilot, and
MAVROS remain unpinned, playing the role of physically separate systems.
All load parameters were calibrated on single Fixed-400 trials and locked
\emph{before} the comparison trials ran.

\subsection{Operating regimes (pre-registered)}

The experimental plan specified a second operating regime --- a reduced
control period --- to evaluate whether conclusions generalize across
timing constraints. The \textbf{primary regime} runs $T = 50$\,ms
(20\,Hz); the \textbf{secondary regime} runs $T = 25$\,ms (40\,Hz) with
the same trajectory, controller code, load profile, and analysis
pipeline. Under Eq.~\eqref{eq:budget} the adaptive budget target scales
with $T$ automatically; the constant conditions, by construction, do not.

\subsection{Trial protocol and metrics}

Each trial receives a \emph{fresh} simulator stack (PX4, Gazebo, MAVROS
restarted per trial) to guarantee independent, uncontaminated initial
conditions; an automated runner enforces identical load-window timing
across trials and a validity guard (circuit completion and steady-state
attainment) with re-runs, so intermittent simulator faults cannot
silently enter the dataset. The primary regime comprises $n = 13$ valid
trials per condition; the secondary, $n = 8$.

\emph{Metrics.} Deadline misses (count of cycles with $c_t > T$); RMS
tracking error, both whole-trajectory and \emph{steady-state-only},
where steady state begins algorithmically at the first run of 10
consecutive cycles with position error below 0.5\,m (no hardcoded time
cutoff); mean $N$; mean call time; state-to-command latency; and
scheduler-stability metrics $\overline{|\Delta N|}$ and
$\mathrm{Var}(N)$, with $\Delta N_t = N_t - N_{t-1}$, to detect
oscillatory scheduling (``chatter''). Scheduler transient behavior is
quantified by response latency (cycles from a per-sample cost spike to a
$\geq$10\% budget reduction) and recovery time. Between-condition
comparisons use Welch's $t$-test and the Mann--Whitney $U$ test on
per-trial statistics, with Holm correction across the reported family of
comparisons and a leave-one-out sensitivity analysis for every pairwise
test (reported wherever the two test types disagree). All headline
deadline-miss results survive Holm correction ($p < 0.005$ adjusted);
all in-regime null results remain null (adjusted $p \approx 1$).

% =========================================================================
\section{Results}
\label{sec:results}

\subsection{Primary regime: 20\,Hz (50\,ms deadline)}
\label{sec:results:20hz}

Table~\ref{tab:20hz} summarizes the primary comparison.

\begin{table}[t]
\caption{Primary regime (20\,Hz, 50\,ms deadline), mean $\pm$ std;
$n{=}13$ per condition ($n{=}8$ for Constant-200, added post-review
under pre-registered interpretation rules). Load window: cores 0--1, 6
threads, 30\,s.}
\label{tab:20hz}
\centering
\begin{tabular}{lcccc}
\toprule
 & Adaptive & Const-200 & Const-330 & Fixed-400 \\
\midrule
Deadline misses      & $0.15 \pm 0.38$  & $0.25 \pm 0.46$ & $0.23 \pm 0.60$  & $18.0 \pm 4.1$ \\
Miss rate            & $0.0\%$          & $0.0\%$         & $0.0\%$          & $1.0\%$ \\
SS-RMS [m]           & $0.131 \pm 0.055$& $0.156 \pm 0.093$ & $0.135 \pm 0.051$& $0.155 \pm 0.071$ \\
Mean $N$             & $365.1$          & $200$           & $330$            & $400$ \\
In-window call [ms]  & $30.0$           & $23.0$          & $32.8$           & $40.6$ \\
$\overline{|\Delta N|}$ & $2.0$         & $0$             & $0$              & $0$ \\
\bottomrule
\end{tabular}
\end{table}

\emph{Adaptive vs.\ Fixed-400.} The adaptive controller essentially
eliminated deadline misses ($0.15 \pm 0.38$ vs.\ $18.0 \pm 4.1$ per
trial; Welch $t = -15.5$, $p < 10^{-3}$; Mann--Whitney $U = 0$,
$p < 10^{-3}$), while steady-state tracking error was statistically
indistinguishable (Welch $p = 0.35$). During the load window the
scheduler held call time at $30.0$\,ms --- exactly its target
$\beta T$ --- by regulating $N$ to $\approx 300$, recovering to
$N_{\max}$ within cycles of load release.
\todo{Figure: 20\,Hz overlay time series (N, call time, error), three
conditions.}

\emph{Adaptive vs.\ the constant budgets.} Under the tested condition,
the predetermined constant budget performed equivalently to adaptation
on both deadline misses ($p = 0.70$) and tracking error ($p = 0.83$;
all in-regime constant comparisons have Holm-adjusted $p \approx 1$): a
well-chosen constant matches adaptation \emph{in the regime it is tuned
for}. To test whether a much smaller budget would also have sufficed, a
Constant-200 arm (near the Pareto elbow) was added after external
review. To avoid post-hoc goalpost-moving, its interpretation rules
were written and committed to version control before any Constant-200
trial ran, binding us to publish either outcome: a Constant-200 failure
in either regime would support the transfer claim, while Constant-200
succeeding everywhere would narrow our claims to the tuning-free
framing alone. The former occurred (Sec.~\ref{sec:results:40hz}). At
20\,Hz:
at 20\,Hz it, too, is indistinguishable from adaptation
($0.25 \pm 0.46$ misses; tracking within noise) --- in this regime,
\emph{every} reasonable constant ties, exactly as the flat Pareto
region and OS-scheduler shielding predict. Sec.~\ref{sec:discussion}
discusses why, and Sec.~\ref{sec:results:40hz} shows what no constant
can do.

\emph{Scheduler validation.} Across adaptive trials the scheduler
responded to per-sample cost spikes within $1.4$ cycles on average and
recovered within $1.6$ cycles; $\overline{|\Delta N|} \approx 2$
indicates no oscillatory chatter --- the variance of $N$ is dominated by
the intended load-window step, not cycle-to-cycle oscillation.
\todo{Figure: $\Delta N$ distribution per condition.}

\subsection{Secondary regime: 40\,Hz (25\,ms deadline)}
\label{sec:results:40hz}

The constant budget tuned at 20\,Hz does not transfer to the tighter
timing regime --- it misses $\sim$30\% of all control cycles --- while
the adaptive scheduler re-regulates to the new budget automatically.
Table~\ref{tab:40hz} summarizes.

\begin{table}[t]
\caption{Secondary regime (40\,Hz, 25\,ms deadline), mean $\pm$ std over
$n{=}8$ trials per condition. Identical trajectory, code, and load
profile as Table~\ref{tab:20hz}.}
\label{tab:40hz}
\centering
\begin{tabular}{lcccc}
\toprule
 & Adaptive & Const-200 & Const-330 & Fixed-400 \\
\midrule
Deadline misses  & $0.75 \pm 0.71$ & $122.3 \pm 6.9$ & $1048.9 \pm 5.2$ & $876.5 \pm 7.6$ \\
Miss rate        & $0.0\%$         & $3.5\%$         & $30.5\%$         & $26.9\%$ \\
SS-RMS [m]       & $0.136 \pm 0.127$ & $0.137 \pm 0.120$ & $0.095 \pm 0.005$ & $0.130 \pm 0.045$ \\
Mean $N$         & $220.5$         & $200$           & $330$            & $400$ \\
Mean call [ms]   & $15.0$          & $16.2$          & $22.0$           & $26.3$ \\
\bottomrule
\end{tabular}
\end{table}

With zero configuration change, the adaptive controller discovered a
different budget for the different deadline --- $N \approx 252$ idle,
$\approx 143$ under load --- pinning call time at $15.0$\,ms
($= \beta T$) and maintaining essentially zero deadline misses
($0.75 \pm 0.71$ per trial; vs.\ every constant: Welch $p < 10^{-3}$,
Mann--Whitney Holm-adjusted $p < 0.005$, leave-one-out robust) at the
full 40\,Hz rate. The failure of the constants scales with their size:
even the elbow-level Constant-200 --- which matched adaptation
perfectly at 20\,Hz --- misses $122.3 \pm 6.9$ cycles per trial
(3.5\%) here, while Constant-330 and Fixed-400 miss 30.5\% and 26.9\%
respectively. Every fixed budget embeds a regime assumption; only the
size of the violation differs. The constant
conditions remained flyable --- Constant-330's steady-state tracking was
in fact the best of the three --- but violated the real-time contract on
roughly every third cycle, degrading the effective control rate to
$\sim$38\,Hz; the cost of a miss in SITL at this margin is the broken
timing contract and rate degradation, not tracking collapse
(Sec.~\ref{sec:discussion}). On steady-state RMS between adaptive and
Fixed-400, the mean-based and rank-based tests disagree (Welch
$p = 0.90$; Mann--Whitney $p = 0.01$), and a leave-one-out sensitivity
analysis resolves the disagreement: the rank result is robust to
removing any single trial (Mann--Whitney $p \in [0.0003, 0.021]$ over
all leave-one-out subsets; Holm-adjusted $p = 0.031$), while Welch is
not ($p \in [0.047, 0.97]$). The per-trial data explain why: adaptive's
seven non-outlier trials (0.081--0.097\,m) all track better than every
Fixed-400 non-outlier trial (0.098--0.126\,m; medians 0.093 vs.\
0.116\,m), and a single adaptive outlier trial (0.451\,m, still a valid
flight) inflates the mean-based comparison. Typical (median) tracking
under adaptation was thus, if anything, \emph{better} at the tighter
deadline; we nevertheless claim only non-inferiority (H2), consistent
with the primary regime.
\todo{Figure: 40\,Hz overlay time series, three conditions.}

\subsection{Safety-fallback validation by fault injection}
\label{sec:results:fault}

Because the budget law's margin prevented any fallback activation under
the comparison conditions, we validated the fallback path directly: in
two additional adaptive trials at 20\,Hz, a synthetic 60\,ms compute
stall was injected inside the timed control path for a 3-cycle window
at $t{=}60$\,s (well clear of the load window), yielding measured call
times of 83--89\,ms against the 50\,ms deadline. In both trials the
sequence unfolded identically: the first stalled cycle registered a
deadline miss and the scheduler immediately cut the budget
($N: 400 \to \approx 255$); the second consecutive miss triggered the
fallback signal exactly at the specified threshold (logged and recorded
in the cycle telemetry); the stall window then expired, the fallback
condition cleared on the next compliant cycle, and the budget recovered
to $N_{\max}$ within at most 10 cycles (0.5\,s). Peak tracking
disturbance during and after the episode was 0.13\,m, and both flights
completed the full circuit normally. The trigger, signal, and recovery
path of the fallback mechanism thus behave as specified under a
worst-case compute transient severe enough that no fair-share load on
this platform could produce it.

\subsection{Equivalent-budget analysis: the Pareto curve predicts flight}
\label{sec:results:budget}

Interpolating the offline Pareto table (Sec.~\ref{sec:pareto}) at each
adaptive trial's mean $N$ predicts the quality cost of adaptation. At
20\,Hz (mean $N = 365.1$), the predicted quality gap versus $N = 400$ is
$-0.0015$ --- 5\% of the table's own noise floor --- with a predicted
per-call compute saving of 7.6\% whole-trial ($\approx$22\% during the
load window). At 40\,Hz (mean $N = 220.5$), the predicted gap is
$+0.0082$, still only $0.27\times$ the noise floor, with predicted
savings of 41.8\% whole-trial ($\approx$63\% in-window). Both
predictions are confirmed by the flight data: no statistically
consistent steady-state tracking difference was observed in either
regime. The offline characterization thus \emph{predicts} closed-loop
outcomes: adaptation operates on the flat region of the quality--compute
curve, which is precisely why halving the budget cost nothing
measurable.

% =========================================================================
\section{Discussion}
\label{sec:discussion}

\emph{What adaptation buys.} The experimental story is: at 20\,Hz,
adaptive $\approx$ tuned constant $>$ untuned constant; at 40\,Hz,
adaptive adapts automatically while the previously tuned constant no
longer satisfies the timing requirements. The value of the adaptive
scheduler is robustness across operating regimes without offline
retuning --- not outperforming the best hand-tuned constant within every
fixed regime. Obtaining Constant-330 required an offline Pareto sweep
plus a stress calibration; the scheduler found the corresponding budget
online, in every regime it encountered, from nothing but its own
measured execution times.

\emph{Why the tuned constant is hard to beat in-regime.} Two measured
platform properties compress the room between ``tuned constant
suffices'' and ``flight fails'' at the 50\,ms deadline. First, the
quality axis: beyond the Pareto elbow, quality differences between
budgets in $[200, 400]$ are within noise --- so no in-regime experiment
can separate budgets by tracking quality. Second, the deadline axis: the
OS scheduler shields a periodic control task from fair-share CPU
contention; across our calibration ladder, per-sample cost inflation
saturated at $\sim$1.4$\times$ regardless of load size, which a 34\%
budget headroom absorbs. We report this as a finding with its own
practical implication: on multi-core companion computers running
general-purpose schedulers, moderate fair-share contention degrades a
periodic controller far less than naive utilization arithmetic suggests
--- and conversely, the regimes that \emph{do} break a tuned constant are
changes in the timing requirement itself, precisely where a ratio law
like Eq.~\eqref{eq:budget} transfers and a constant cannot.

\emph{On the cost of a miss.} In our SITL experiments, missing the
deadline on a third of cycles degraded the effective control rate but
not steady-state tracking --- the velocity-setpoint interface is
forgiving at these margins. We deliberately do not claim that misses
caused tracking failure; the claim is that the adaptive controller
maintains the real-time contract that autopilot stacks assume, and that
weakly-hard analyses~\cite{maggio2020control} show cannot be violated
arbitrarily. Platforms with tighter inner loops, attitude-level
interfaces, or hard real-time certification requirements pay a higher
price per miss.

\emph{Composability.} The scheduler touches MPPI only through the
per-call sample count and measured duration. It is therefore orthogonal
to, and composable with, distribution-adaptation
methods~\cite{asmar2023mpopis,mohamed2022logmppi}, robustness
layers~\cite{gandhi2021robust,yin2023shield}, and GPU
parallelization~\cite{minarik2024mppi,enrico2025comparison} --- for a
GPU implementation the per-sample cost model would need to account for
batch effects, which we leave to future work.

% =========================================================================
\section{Threats to Validity}
\label{sec:threats}

\emph{Internal.} MPPI is stochastic and Gazebo physics is variable;
run-to-run variance is addressed by repeated trials with a fresh
simulator stack per trial, algorithmic (not hand-picked) steady-state
segmentation, and rank-based tests alongside parametric ones. Scheduler
timing is subject to OS jitter on the test machine, independent of the
algorithm.

\emph{External.} Results are software-in-the-loop on a single quadrotor
model, task (square circuit and station-keeping), and desktop-class CPU;
the prediction model is a kinematic single-integrator over velocity
commands. The 40\,Hz regime represents tighter real-time constraints of
the kind expected on lower-performance embedded processors, but we make
no equivalence claim to any specific embedded platform without
benchmarking it; hardware-in-the-loop validation is planned future work.

\emph{Construct.} Busy-wait CPU load approximates contention; our own
calibration shows the OS shields periodic tasks from it, and a
memory-bandwidth variant of the load generator produced similar
saturation on this compute-bound controller. The two-core pin emulates
an embedded compute envelope rather than reproducing any specific
device. The state-to-command latency we report upper-bounds middleware
overhead, as it includes the asynchronous gap between pose arrival and
the control timer.

% =========================================================================
\section{Conclusion}
\label{sec:conclusion}

We treated MPPI's sample count as a runtime-adaptive control variable
governed by a measured-execution-time budget law with a deadline
objective and safety fallback, and validated it in staged, pre-registered
SITL experiments against both the conventional fixed-maximum budget and a
predetermined equivalent-budget constant, across two timing regimes. The
adaptive controller essentially eliminated deadline misses under CPU
contention where the conventional baseline missed consistently; matched a
carefully tuned constant in the regime that constant was tuned for; and,
unlike the constant, transferred to a 2$\times$ tighter deadline with
zero retuning, re-discovering an appropriate budget online while the
tuned constant missed $\sim$30\% of cycles. An offline quality--compute
characterization predicted these closed-loop outcomes, because adaptation
operates on the flat region of MPPI's quality--compute curve.

Future work: a Pareto-informed scheduler that additionally enforces an
empirical quality floor (the current law is deadline-only);
hardware-in-the-loop and onboard-GPU deployment, where per-sample cost
models must account for batch effects; and composition with
distribution-adaptive MPPI variants.

% =========================================================================
\bibliographystyle{IEEEtran}
\bibliography{references}

\end{document}
